\documentclass{article}

% Language setting
% Replace `english' with e.g. `spanish' to change the document language
\usepackage[english]{babel}

% Set page size and margins
% Replace `letterpaper' with `a4paper' for UK/EU standard size
\usepackage[letterpaper,top=2cm,bottom=2cm,left=3cm,right=3cm,marginparwidth=1.75cm]{geometry}

% Useful packages
\usepackage{amsmath}
\usepackage{amsthm}
\usepackage{dsfont}
\usepackage{graphicx}
\usepackage{forest}
\usepackage[colorlinks=true, allcolors=blue]{hyperref}
\usepackage[mathscr]{euscript}

\usepackage{physics}

%NECESRY
\usepackage{principia}
\theoremstyle{definition}
\newtheorem{definition}{Definition}[section]
\newtheorem{remark}{Remark}[section]
\usepackage{bxcjkjatype}

\title{Ground Truth Ideas}
\author{You}

\begin{document}
\maketitle

\section{Assumptions}
Suppose there are mutually independent features $x_1^I,\ldots,x_d^I$ and features $x_{d+1}^C,\ldots,x_D^C$ that are each (possibly nonlinear) functions of the mutually independent features. We have a data sample $\mathcal{D} = \{\mathbf{x}^{(n)}\}_{n=1}^N$
% That is, there are functions $f_{d+1},\ldots,f_D$ such that:
% \begin{gather*}
%     x_{d+1}^C = f_{d+1}(x_1^I,\ldots,x_d^I) \\
%     x_{d+2}^C = f_{d+2}(x_1^T,\ldots,x_d^I) \\
%     \ldots \\
%     x_{D}^C = f_D(x_1^T,\ldots,x_d^I) \\ 
% \end{gather*}

Let the ground truth model be a (nonlinear) function $g$ of the features:
\begin{equation*}
    \mathcal{G}(x) = g(x_1^I,\ldots,x_d^I,x_{d+1}^C,\ldots,x_D^C)
\end{equation*}

We can write this as a composite function of only the independent variables:
\begin{equation*}
    \widetilde{g}(f_{d+1}(x_1^I,\ldots,x_d^I),\ldots,f_D(x_1^I,\ldots,x_D^I))
\end{equation*}

\section{Estimand In Practice}
\begin{equation*}
    \text{importance of } x_j = \frac{1}{N} \sum_{n=1}^N \mid \text{SHAP}_\mathcal{M}(x_j^{(n)}) \mid 
\end{equation*}
where $\mathcal{M}$ is some XGBoost model. 

\section{Ground Truth Possibilities}
    \subsection{Total Derivative}
        \textbf{Meaning}: $\dv{\widetilde{g}}{x_1^I}$ is the influence of $x_I^I$ on $g$ considering every appearance of $x_1^I$ in all arguments of $\widetilde{g}$

        \bigskip

        % https://math.stackexchange.com/questions/4131291/is-the-total-derivative-defined-for-a-function-of-multiple-independent-variables
        \textbf{Definition}:
        \begin{gather*}
            \dv{\widetilde{g}}{x_1^I} = \sum_{d'=d+1}^D \pdv{\widetilde{g}}{f_{d'}} \pdv{f_{d'}}{x_1^I} \\
            \ldots \\
            \dv{\widetilde{g}}{x_d^I} = \sum_{d'=d+1}^D \pdv{\widetilde{g}}{f_{d'}} \pdv{f_{d'}}{x_d^I} \\
        \end{gather*}
    
        % For $d'=1,\ldots,d$:
        % \begin{align*}
        %     \dv{g(x_1^I,\ldots,x_d^I,x_{d+1}^C,\ldots,x_D^C)}{x_{d'}} = \pdv{g}{x_{d'}} + \sum_{d''=d+1}^D \pdv{g}{x_{d''}} \pdv{x_{d''}}{x_d'}
        % \end{align*}

        % \bigskip

        % For $d'=d+1,\ldots,D$ there is no total derivative $\dv{g}{x_{d'}}$

    \subsection{Marginal Effect}
        \subsubsection{Average Marginal Effects (AME)}
        \textbf{Meaning}: Estimate of the expected marginal effect w.r.t. the data distribution, using a sample of data $\mathcal{D} = \{\mathbf{x}^{(n)}\}_{n=1}^N$ for the estimate. \textbf{Given the prevalent use of mean absolute SHAP values, I think this version of marginal effect is the closest marginal effect-based analogue of a ground truth.}

        \bigskip

        \textbf{Definition}: 
        %Let $\text{fME}_{\mathbf{x}^{(n)},h_j} = g(x_1,\ldots,x_j + h_j,\ldots,X_D) - g(\mathbf{x}^{(n)})$ 
        \begin{equation*}
            \text{AME}_{\mathcal{D},h_j} = \frac{1}{n}\sum_{n=1}^N \left[ g(x_1^{(n)},\ldots,x_j^{(n)} + h_j,\ldots,x_D^{(n)}) - g(\mathbf{x}^{(n)}) \right]
        \end{equation*}
        where $h_j$ is some perturbation.
    
        % \begin{equation*}
        %     \mathds{E}_X\left[\pdv{g}{x_{d'}}\right]
        % \end{equation*}

        \subsubsection{Marginal Effect at Means (MEM)}
        \textbf{Meaning}: ME evaluated at the expectation of X

        \bigskip

        \textbf{Definition}:
        \begin{equation*}
            \text{MEM}_{\mathcal{D},h_j} = g(\bar{x}_1,\ldots,\bar{x}_j + h_j,\ldots,\bar{x}_D) - g(\bar{\mathbf{x}})
        \end{equation*}

    
        % \begin{equation*}
        %     \pdv{g}{x_{d'}}\mid_{x=\bar{x}}
        % \end{equation*}

        \subsubsection{Marginal Effect for Categorical Feature}
        \textbf{Meaning}:         For each observation, the observed categorical feature value is set to the reference category and we record the change in prediction when changing it to every other category. Given $k$ categories, this results in $k-1$ MEs for each observation. 

        \bigskip

        \textbf{Definition}: 

    



        \begin{equation*}
            g(c',X_{-d}) - g(c,X_{-d})
        \end{equation*}
\end{document}